\section{Gabarito e resoluções comentadas — Controle Externo}

\begin{solucao}{06.01}\textbf{CERTO.} É o desenho do art. 71 combinado com o art. 70 da Constituição Federal.\end{solucao}
\begin{solucao}{06.02}\textbf{ERRADO.} O Tribunal auxilia o Legislativo, mas não se subordina hierarquicamente a ele.\end{solucao}
\begin{solucao}{06.03}\textbf{CERTO.} Legalidade, legitimidade e economicidade integram expressamente a fiscalização constitucional.\end{solucao}
\begin{solucao}{06.04}\textbf{ERRADO.} O TCU emite parecer prévio; o julgamento político compete ao Congresso Nacional.\end{solucao}
\begin{solucao}{06.05}\textbf{CERTO.} É competência constitucional própria do Tribunal.\end{solucao}
\begin{solucao}{06.06}\textbf{ERRADO.} Irregularidade sem dano pode gerar outros efeitos, mas débito exige dano e pressupostos de responsabilização.\end{solucao}
\begin{solucao}{06.07}\textbf{CERTO.} O sistema constitucional prevê integração e apoio, não substituição total.\end{solucao}
\begin{solucao}{06.08}\textbf{ERRADO.} Denúncia é notícia/peça de provocação; os fatos precisam ser instruídos e comprovados.\end{solucao}
\begin{solucao}{06.09}\textbf{CERTO.} A cautelar preserva utilidade e previne dano; não declara responsabilidade definitiva.\end{solucao}
\begin{solucao}{06.10}\textbf{ERRADO.} Monitoramento deve verificar cumprimento material e, quando relevante, efetividade.\end{solucao}
\begin{solucao}{06.11}\textbf{A.} O Tribunal emite parecer prévio e o Legislativo julga as contas de governo.\end{solucao}
\begin{solucao}{06.12}\textbf{A.} Administradores e responsáveis submetem-se ao julgamento técnico do Tribunal de Contas, conforme a competência aplicável.\end{solucao}
\begin{solucao}{06.13}\textbf{CERTO.} Contas de governo e contas de gestão têm objetos, responsáveis e efeitos distintos.\end{solucao}
\begin{solucao}{06.14}\textbf{B.} Economicidade considera custo global, qualidade e alternativas, não o preço inicial isolado.\end{solucao}
\begin{solucao}{06.15}\textbf{A.} A legitimidade permite examinar aderência à finalidade e ao interesse público.\end{solucao}
\begin{solucao}{06.16}\textbf{CERTO.} Auditoria operacional usa critérios de desempenho, eficiência e efetividade.\end{solucao}
\begin{solucao}{06.17}\textbf{B.} Falha formal não prova dano. Débito exige quantificação e nexo de responsabilidade.\end{solucao}
\begin{solucao}{06.18}\textbf{A.} Responsabilização exige conexão entre conduta e resultado dentro do regime jurídico aplicável.\end{solucao}
\begin{solucao}{06.19}\textbf{CERTO.} Solidariedade depende de fundamento jurídico e demonstração dos pressupostos.\end{solucao}
\begin{solucao}{06.20}\textbf{B.} Cautelar serve para evitar dano e preservar a eficácia da decisão final.\end{solucao}
\begin{solucao}{06.21}\textbf{CERTO.} A decisão urgente não substitui o exame de mérito e o devido processo posterior.\end{solucao}
\begin{solucao}{06.22}\textbf{A.} O controle interno possui funções próprias e de apoio ao externo, nos termos constitucionais.\end{solucao}
\begin{solucao}{06.23}\textbf{B.} Representação deve passar por admissibilidade e instrução; documentos são avaliados, não presumidos conclusivos.\end{solucao}
\begin{solucao}{06.24}\textbf{CERTO.} A fiscalização pode ser concomitante sem assumir competência gerencial.\end{solucao}
\begin{solucao}{06.25}\textbf{B.} Evidência de revogação e do processo preventivo demonstra cumprimento real, não apenas declaração.\end{solucao}
\begin{solucao}{06.26}\textbf{C.} I e III corretas; II é falsa.\end{solucao}
\begin{solucao}{06.27}\textbf{CERTO.} O Tribunal possui competências constitucionais próprias, inclusive julgamento e sanção.\end{solucao}
\begin{solucao}{06.28}\textbf{A.} O parecer prévio subsidia o julgamento das contas de governo pelo Legislativo.\end{solucao}
\begin{solucao}{06.29}\textbf{B.} A competência de registro opera segundo Constituição, lei e jurisprudência aplicáveis.\end{solucao}
\begin{solucao}{06.30}\textbf{CERTO.} Trata-se de controle externo administrativo especializado, não jurisdição judicial.\end{solucao}
\begin{solucao}{06.31}\textbf{B.} A TCE é mecanismo excepcional de apuração de dano/responsabilidade quando o fluxo ordinário é insuficiente, conforme a norma aplicável.\end{solucao}
\begin{solucao}{06.32}\textbf{B.} Sanção sem defesa viola garantias processuais fundamentais.\end{solucao}
\begin{solucao}{06.33}\textbf{CERTO.} Cautelares podem comportar contraditório diferido em situações urgentes, segundo o regime jurídico.\end{solucao}
\begin{solucao}{06.34}\textbf{B.} Multa exige base legal, competência, responsabilidade e devido processo.\end{solucao}
\begin{solucao}{06.35}\textbf{B.} Ordem vinculada ao cumprimento de dever jurídico é determinação.\end{solucao}
\begin{solucao}{06.36}\textbf{CERTO.} Recomendações podem focar resultado, evitando prescrição tecnológica desnecessária.\end{solucao}
\begin{solucao}{06.37}\textbf{B.} Prescrição deve ser examinada conforme a pretensão e o regime vigente, não por fórmulas absolutas.\end{solucao}
\begin{solucao}{06.38}\textbf{B.} O caso exige avaliar desempenho e resultados, além da conformidade formal.\end{solucao}
\begin{solucao}{06.39}\textbf{CERTO.} TI é objeto legítimo de controle quando afeta gestão, recursos, serviços e riscos públicos.\end{solucao}
\begin{solucao}{06.40}\textbf{B.} Accountability envolve explicar, justificar e responder pela gestão e seus resultados.\end{solucao}
\begin{solucao}{06.41}\textbf{CERTO.} A Constituição autoriza exame de economicidade independentemente da formalidade aparente.\end{solucao}
\begin{solucao}{06.42}\textbf{B.} Dano sem nexo com o agente não basta para imputação pessoal.\end{solucao}
\begin{solucao}{06.43}\textbf{B.} Escolher especificações e autorizar decisões gerenciais compromete a independência do controle.\end{solucao}
\begin{solucao}{06.44}\textbf{CERTO.} O controle deve prevenir sem coadministrar.\end{solucao}
\begin{solucao}{06.45}\textbf{A.} A cautelar é provisória e instrumental; a sanção é consequência definitiva de responsabilidade apurada.\end{solucao}
\begin{solucao}{06.46}\textbf{CERTO.} Aprovação pelo controle interno não vincula o controle externo.\end{solucao}
\begin{solucao}{06.47}\textbf{B.} Requisitos processuais de admissibilidade e prova continuam válidos mesmo diante de alegação grave.\end{solucao}
\begin{solucao}{06.48}\textbf{B.} Criar política é implementação documental; restore falhando demonstra falta de efetividade.\end{solucao}
\begin{solucao}{06.49}\textbf{CERTO.} O auditor deve evitar assumir a escolha gerencial quando há alternativas válidas.\end{solucao}
\begin{solucao}{06.50}\textbf{B.} O achado correto delimita ausência de evidência, risco e procedimentos necessários, sem acusação prematura.\end{solucao}
